\section{Introduction}\label{sec:intro}

After the 2022 attacks on SIDH~\cite{CD23,MMPPW23,Rob23}, isogeny-based public-key
encryption needs a trapdoor that does not reveal enough torsion information to
let the attacker run those same attacks. Two families
survive. In the first, torsion images are masked so that the attacker cannot
recover them exactly, as in M-SIDH~\cite{FMP23} and FESTA~\cite{BMP23}; in the
second, the scheme avoids torsion images altogether. FESTA decrypts by running
the torsion attack itself with the help of a secret: the legitimate receiver
knows a matrix that commutes with the masks.

Castryck and Vercauteren~\cite{CV23} observed that a torsion basis which is an
eigenbasis of a \emph{known} endomorphism of small degree makes masked-torsion
instances easy: the masks act diagonally on the eigenvectors, so the
endomorphism's action on the masked points is known exactly, and a lollipop
endomorphism can be decoded with Kani's lemma. They presented this as an
attack, and as a backdoor available to whoever chose the basis.

\paragraph{This work.} Eigen-S\'eta turns that observation into a trapdoor. The
public key is a curve $\EA$ with an eigenbasis of an endomorphism
$\tau\in\End(\EA)$ whose degree $2f^2$ is \emph{secret}. The sender masks the
images of the eigenbasis under a $3^F$-isogeny $\varphi$ with a random diagonal
matrix. The receiver, who knows $f$, knows how
$\theta=\varphi\tau\hat\varphi$ acts on the masked points and recovers $\varphi$
from one $(2^E,2^E)$-isogeny between products of elliptic curves. Without $f$,
the ciphertext is a FESTA-type masked-torsion instance, and the public key is
an isogeny instance of unknown degree from a known curve.

An earlier variant~\cite{EigenSetaNote} took $\tau=\varepsilon\omega\hat\varepsilon$
with $\varepsilon$ of \emph{public} degree $5^k$. That variant is insecure:
its decryption key is a function of the public degree alone
(\S\ref{sec:walk}). Version~2, specified here, draws $f$ from an interval of
about $2^{277}$, $2^{408}$ or $2^{535}$ values, depending on the level, and keeps it secret.

\paragraph{Contributions.}
\begin{itemize}
\item A complete specification of the KEM (\S\ref{sec:scheme}) with the
  Fujisaki--Okamoto transform with implicit rejection~\cite{HHK17}.
\item Proofs that decryption is correct (\S\ref{sec:correctness}). These
  include a proposition that the kernel recovered during decryption is always
  cyclic, which settles a corner case reported for the earlier toy prototype.
\item Parameter sets for $\lambda=128,192,256$ (\S\ref{sec:params}) and a
  heuristic analysis of key recovery, message recovery, quantum attacks,
  chosen-ciphertext attacks and decryption failures (\S\ref{sec:security}).
  The concrete figures are summarised in Table~\ref{tab:attacks}.
\item A Rust implementation with constant-time decapsulation, checked with an
  instruction-count probe and a statistical timing test. The probe found one
  compiler-introduced branch, which we fixed (\S\ref{sec:impl}).
\item Measured sizes, timings and correctness rates at all parameter sets, and
  constant-time evidence at Toy16 and L1 (\S\ref{sec:results}), each
  traceable to a committed record.
\end{itemize}

\paragraph{What this paper does not claim.} Security is a conjecture backed by
the attacks we considered; no third party has analysed the scheme. The
implementation is research code. Its decapsulation is slow, mainly because
the dimension-2 chain is computed in Mumford coordinates; a theta-coordinate
chain~\cite{DMPR23} is future work. Key generation is not constant time.

\paragraph{Evidence conventions.} Each measured number carries the path of the
record it comes from, relative to the root of the repository
(\art{crypto:}). Derived quantities are marked as such.
